\documentclass[11pt,a4paper]{article}
\usepackage[top=2.3cm,bottom=2.3cm,left=2.2cm,right=2.2cm]{geometry}

\usepackage{amsmath,amssymb}
\usepackage{fontspec}
\setmainfont{Liberation Serif}
\setsansfont{Liberation Sans}
\setmonofont{DejaVu Sans Mono}[Scale=0.84]

\usepackage{xcolor}
\definecolor{brandblue}{RGB}{20, 55, 105}
\definecolor{accentblue}{RGB}{35, 95, 165}
\definecolor{codebg}{RGB}{248, 249, 251}
\definecolor{codeframe}{RGB}{215, 222, 230}
\definecolor{javakeyword}{RGB}{0, 45, 170}
\definecolor{javastring}{RGB}{5, 120, 20}
\definecolor{javacomment}{RGB}{120, 120, 120}
\definecolor{javanumber}{RGB}{20, 75, 220}

\definecolor{noteborder}{RGB}{30, 110, 65}
\definecolor{notebg}{RGB}{242, 249, 245}
\definecolor{warnborder}{RGB}{185, 75, 10}
\definecolor{warnbg}{RGB}{254, 248, 240}
\definecolor{tipborder}{RGB}{30, 95, 160}
\definecolor{tipbg}{RGB}{242, 247, 254}

\usepackage{fancyhdr}
\usepackage{lastpage}
\pagestyle{fancy}
\fancyhf{}
\lhead{\parbox[t]{0.58\textwidth}{\raggedright\small\sffamily\color{brandblue}\textbf{T16 --- Java Collections Framework (JCF)}}}
\rhead{\small\sffamily\color{gray}Programmazione II -- UniVR (2025/2026)}
\cfoot{\small\sffamily Pagina \thepage\ di \pageref{LastPage}}
\renewcommand{\headrulewidth}{0.5pt}
\renewcommand{\headrule}{\hbox to\headwidth{\color{brandblue!70!black}\leaders\hrule height \headrulewidth\hfill}}

\usepackage{titlesec}
\titleformat{\section}{\Large\bfseries\sffamily\color{brandblue}}{\thesection}{0.8em}{}[\titlerule]
\titleformat{\subsection}{\large\bfseries\sffamily\color{accentblue}}{\thesubsection}{0.8em}{}
\titleformat{\subsubsection}{\normalsize\bfseries\sffamily\color{brandblue!85!black}}{\thesubsubsection}{0.8em}{}

\usepackage{needspace}
\usepackage{fancyvrb}
\DefineVerbatimEnvironment{asciibox}{Verbatim}{
  fontsize=\fontsize{7.5pt}{9.2pt}\selectfont,
  frame=single,
  rulecolor=\color{codeframe},
  fillcolor=\color{codebg},
  framesep=2.5mm
}
\usepackage{listings}
\lstset{
  backgroundcolor=\color{codebg},
  frame=single,
  rulecolor=\color{codeframe},
  basicstyle=\ttfamily\footnotesize,
  keywordstyle=\color{javakeyword}\bfseries,
  stringstyle=\color{javastring},
  commentstyle=\color{javacomment}\itshape,
  numberstyle=\tiny\color{gray},
  numbers=left,
  numbersep=7pt,
  stepnumber=1,
  breaklines=true,
  breakatwhitespace=true,
  showstringspaces=false,
  tabsize=4,
  xleftmargin=12pt,
  framexleftmargin=12pt
}

\newcommand{\passthrough}[1]{#1}
\providecommand{\tightlist}{\setlength{\itemsep}{0pt}\setlength{\parskip}{0pt}}

\usepackage{tcolorbox}
\tcbuselibrary{skins,breakable}
\newtcolorbox{studynote}[1][Nota]{
  enhanced,
  breakable,
  colback=notebg,
  colframe=noteborder,
  fonttitle=\bfseries\sffamily\small,
  coltitle=white,
  title={#1},
  arc=2mm,
  boxrule=0.8pt,
  left=9pt, right=9pt, top=7pt, bottom=7pt
}

\newtcolorbox{studywarning}[1][Attenzione]{
  enhanced,
  breakable,
  colback=warnbg,
  colframe=warnborder,
  fonttitle=\bfseries\sffamily\small,
  coltitle=white,
  title={#1},
  arc=2mm,
  boxrule=0.8pt,
  left=9pt, right=9pt, top=7pt, bottom=7pt
}

\newtcolorbox{studytip}[1][Suggerimento]{
  enhanced,
  breakable,
  colback=tipbg,
  colframe=tipborder,
  fonttitle=\bfseries\sffamily\small,
  coltitle=white,
  title={#1},
  arc=2mm,
  boxrule=0.8pt,
  left=9pt, right=9pt, top=7pt, bottom=7pt
}

\usepackage{booktabs}
\usepackage{longtable}
\usepackage{array}
\usepackage{calc}

\usepackage{hyperref}
\hypersetup{
  colorlinks=true,
  linkcolor=accentblue,
  urlcolor=accentblue,
  citecolor=brandblue,
  pdfborder={0 0 0}
}

\title{\Huge\bfseries\sffamily\color{brandblue} T16 --- Java Collections Framework (JCF) \\[0.3em] \large\sffamily Gerarchia Collection, List, Set, Map, Queue; Prestazioni e Scelta delle Strutture Dati}
\author{\textbf{Docente:} Prof. Michele Pasqua \quad---\quad \textbf{Appunti Revisione:} Wally}
\date{A.A. 2025/2026 -- Universit\`a degli Studi di Verona}

\begin{document}
\maketitle
\thispagestyle{fancy}
\vspace{-1.2em}
\noindent\textcolor{brandblue}{\rule{\textwidth}{0.8pt}}
\vspace{1.2em}

\begin{center}\rule{0.5\linewidth}{0.5pt}\end{center}

\subsection{1. Analisi Teorica Approfondita (Slide
T16)}\label{analisi-teorica-approfondita-slide-t16}

La lezione 16 è dedicata al \textbf{Java Collections Framework (JCF)}
(\passthrough{\lstinline!java.util!}), una delle architetture cardine
del linguaggio Java: un insieme unificato e standardizzato di
interfacce, implementazioni concrete e algoritmi polimorfici per la
gestione di collezioni di oggetti.

\begin{center}\rule{0.5\linewidth}{0.5pt}\end{center}

\subsubsection{1.1 Architettura e Gerarchia delle Collezioni (Slide
3-4)}\label{architettura-e-gerarchia-delle-collezioni-slide-3-4}

Tutte le collezioni di elementi singoli discendono dall'interfaccia
radice \passthrough{\lstinline!Iterable<E>!}:

\needspace{7cm}
\begin{asciibox}
┌───────────────────┐
│    Iterable<E>    │
└─────────▲─────────┘
          │
┌─────────┴─────────┐
│   Collection<E>   │
└─────────▲─────────┘
          │
┌─────────┴───────────────┬─────────────────────────┐
│                         │                         │
┌─────────┴─────────┐     ┌─────────┴─────────┐     ┌─────────┴─────────┐
│      List<E>      │     │      Set<E>       │     │     Queue<E>      │
└─────────▲─────────┘     └─────────▲─────────┘     └─────────▲─────────┘
          │                         │                         │
   ┌──────┴──────┐           ┌──────┴──────┐           ┌──────┴──────┐
   │             │           │             │           │             │
┌──┴────┐   ┌────┴────┐ ┌────┴────┐   ┌────┴────┐ ┌────┴────┐   ┌────┴────┐
│Array- │   │ Linked- │ │ HashSet │   │ TreeSet │ │Priority-│   │  Deque  │
│ List  │   │  List   │ │         │   │(Sorted) │ │  Queue  │   │(Double) │
└───────┘   └────┬────┘ └─────────┘   └─────────┘ └─────────┘   └────┬────┘
                 │                                                   │
                 └───────────────────┬───────────────────────────────┘
                                     │ (LinkedList implementa
                                     ▼  sia List che Deque/Queue)
\end{asciibox}

\begin{studynote}[Nota di Studio] L'interfaccia \passthrough{\lstinline!Map<K, V>!} (dizionari
chiave-valore) \textbf{NON estende
\passthrough{\lstinline!Collection<E>!}}, poiché memorizza coppie
\((K, V)\) anziché singoli elementi. Tuttavia, fa parte integrante a
pieno titolo del \emph{Java Collections Framework}.
\end{studynote}

\begin{center}\rule{0.5\linewidth}{0.5pt}\end{center}

\subsubsection{\texorpdfstring{1.2 L'Interfaccia Radice
\texttt{Collection\textless{}E\textgreater{}} (Slide
6-7)}{1.2 L'Interfaccia Radice Collection\textless E\textgreater{} (Slide 6-7)}}\label{linterfaccia-radice-collectione-slide-6-7}

Definisce il contratto comune a tutti i contenitori di elementi:

\begin{lstlisting}[language=Java]
public interface Collection<E> extends Iterable<E> {
    int size();
    boolean isEmpty();
    boolean contains(Object o);                    // Ricerca basata su equals()
    boolean containsAll(Collection<?> c);
    boolean add(E e);                              // Aggiunge un elemento (restituisce false se rifiutato)
    boolean addAll(Collection<? extends E> c);     // Wildcard covariante!
    boolean remove(Object o);
    boolean removeAll(Collection<?> c);
    void clear();
    Object[] toArray();
    <T> T[] toArray(T[] a);                        // es. set.toArray(new Person[0])
    Iterator<E> iterator();                        // Ereditato da Iterable<E>
}
\end{lstlisting}

\paragraph{Interscambiabilità tra Collezioni (Slide
7)}\label{interscambiabilituxe0-tra-collezioni-slide-7}

Tutte le implementazioni forniscono per convenzione un costruttore che
accetta un'altra \passthrough{\lstinline!Collection<? extends E>!},
permettendo conversioni istantanee:

\begin{lstlisting}[language=Java]
Collection<Person> list = new LinkedList<>();
list.add(new Person("Joe"));
list.add(new Person("Sam"));

// Conversione da List a Set (elimina al volo eventuali duplicati):
Collection<Person> set = new HashSet<>(list);

// Esportazione in Array tipizzato:
Person[] arrayP = set.toArray(new Person[0]); // Idioma standard Java
\end{lstlisting}

\begin{center}\rule{0.5\linewidth}{0.5pt}\end{center}

\subsubsection{\texorpdfstring{1.3 L'Interfaccia
\texttt{List\textless{}E\textgreater{}} e le sue Implementazioni (Slide
8-9)}{1.3 L'Interfaccia List\textless E\textgreater{} e le sue Implementazioni (Slide 8-9)}}\label{linterfaccia-liste-e-le-sue-implementazioni-slide-8-9}

Una \textbf{Lista} è una sequenza ordinata di elementi che: 1.
\textbf{Preserva l'ordine di inserimento}. 2. \textbf{Accetta elementi
duplicati} (possono esistere più elementi \(e_1, e_2\) tali che
\passthrough{\lstinline!e1.equals(e2)!} sia vero). 3. \textbf{Fornisce
accesso posizionale indicizzato} tramite indice intero
\([0 \dots \text{size}-1]\): *
\passthrough{\lstinline!E get(int index)!}: legge l'elemento all'indice
specificato. * \passthrough{\lstinline!E set(int index, E element)!}:
rimpiazza l'elemento alla posizione data. *
\passthrough{\lstinline!void add(int index, E element)!}: inserisce
shiftando a destra gli elementi successivi. *
\passthrough{\lstinline!E remove(int index)!}: elimina shiftando a
sinistra. * \passthrough{\lstinline!int indexOf(Object o)!}: restituisce
l'indice della prima occorrenza (o \(-1\)). *
\passthrough{\lstinline!List<E> subList(int fromIndex, int toIndex)!}:
vista sulla porzione \([fromIndex, toIndex)\).

\paragraph{\texorpdfstring{Confronto delle Implementazioni:
\texttt{ArrayList} vs \texttt{LinkedList} (Slide 5 e
9)}{Confronto delle Implementazioni: ArrayList vs LinkedList (Slide 5 e 9)}}\label{confronto-delle-implementazioni-arraylist-vs-linkedlist-slide-5-e-9}

\begin{longtable}[]{@{}
  >{\raggedright\arraybackslash}p{(\linewidth - 4\tabcolsep) * \real{0.3333}}
  >{\raggedright\arraybackslash}p{(\linewidth - 4\tabcolsep) * \real{0.3333}}
  >{\raggedright\arraybackslash}p{(\linewidth - 4\tabcolsep) * \real{0.3333}}@{}}
\toprule\noalign{}
\begin{minipage}[b]{\linewidth}\raggedright
Operazione
\end{minipage} & \begin{minipage}[b]{\linewidth}\raggedright
\passthrough{\lstinline!ArrayList<E>!} (Array ridimensionabile)
\end{minipage} & \begin{minipage}[b]{\linewidth}\raggedright
\passthrough{\lstinline!LinkedList<E>!} (Lista doppiamente concatenata)
\end{minipage} \\
\midrule\noalign{}
\endhead
\bottomrule\noalign{}
\endlastfoot
\textbf{Accesso per indice (\passthrough{\lstinline!get(i)!})} &
\textbf{\(O(1)\)} (accesso diretto in memoria contigua) &
\textbf{\(O(n)\)} (deve scorrere i nodi dal capo più vicino) \\
\textbf{Inserimento in coda (\passthrough{\lstinline!add(e)!})} &
\textbf{\(O(1)\)} ammortizzato & \textbf{\(O(1)\)} \\
\textbf{Inserimento in testa (\passthrough{\lstinline!add(0, e)!})} &
\textbf{\(O(n)\)} (deve copiare e shiftare l'array) & \textbf{\(O(1)\)}
(cambio puntatori dei nodi) \\
\textbf{Cancellazione (\passthrough{\lstinline!remove(i)!})} &
\textbf{\(O(n)\)} (shift dei successivi) & \textbf{\(O(1)\)} se si ha il
riferimento al nodo, \(O(n)\) per cercarlo \\
\textbf{Overhead di memoria} & Minimo (array contiguo con eventuale
\emph{capacity} vuota) & Alto (ogni nodo crea un oggetto con puntatori
\passthrough{\lstinline!prev!} e \passthrough{\lstinline!next!}) \\
\textbf{Funzionalità extra} &
\passthrough{\lstinline!ensureCapacity(int minCapacity)!} & Implementa
anche \passthrough{\lstinline!Queue!} e \passthrough{\lstinline!Deque!}:
\passthrough{\lstinline!addFirst()!},
\passthrough{\lstinline!addLast()!}, \passthrough{\lstinline!pop()!},
\passthrough{\lstinline!push()!} \\
\end{longtable}

\begin{center}\rule{0.5\linewidth}{0.5pt}\end{center}

\subsubsection{\texorpdfstring{1.4 L'Interfaccia
\texttt{Set\textless{}E\textgreater{}} e le sue Implementazioni (Slide
10-16)}{1.4 L'Interfaccia Set\textless E\textgreater{} e le sue Implementazioni (Slide 10-16)}}\label{linterfaccia-sete-e-le-sue-implementazioni-slide-10-16}

Un \textbf{Insieme (Set)} modella il concetto matematico di insieme: 1.
\textbf{Nessun duplicato consentito}: non possono mai coesistere due
elementi tali che \passthrough{\lstinline!e1.equals(e2)!} sia vero. 2.
\textbf{Nessun accesso posizionale}: non esistono indici, non si può
fare \passthrough{\lstinline!get(i)!}. 3. L'interfaccia
\passthrough{\lstinline!Set<E>!} non aggiunge nuovi metodi rispetto a
\passthrough{\lstinline!Collection<E>!}, ma ne restringe il contratto
semantico.

\paragraph{\texorpdfstring{1. \texttt{HashSet\textless{}E\textgreater{}}
(Slide
11-14)}{1. HashSet\textless E\textgreater{} (Slide 11-14)}}\label{hashsete-slide-11-14}

\begin{itemize}
\tightlist
\item
  \textbf{Struttura interna}: Mantiene internamente una tabella hash
  (\passthrough{\lstinline!HashMap<E, Object>!}).
\item
  \textbf{Complessità}: Ricerca (\passthrough{\lstinline!contains!}),
  inserimento (\passthrough{\lstinline!add!}) e rimozione
  (\passthrough{\lstinline!remove!}) avvengono in \textbf{tempo costante
  medio \(O(1)\)}.
\item
  \textbf{Ordine}: \textbf{Nessun ordine garantito}; l'ordine di
  scansione dell'iteratore può cambiare nel tempo se la tabella viene
  riallocata (\emph{rehash}).
\item
  \textbf{Vincolo Fondamentale (Slide 12-14)}: L'univocità si basa su
  \textbf{\passthrough{\lstinline!equals()!} e
  \passthrough{\lstinline!hashCode()!}}.

  \begin{itemize}
  \tightlist
  \item
    \textbf{La Regola Aurea}: \emph{``Se fai l'override di
    \passthrough{\lstinline!equals()!}, DEVI fare l'override anche di
    \passthrough{\lstinline!hashCode()!}, e viceversa''}.
  \item
    Proprietà vincolante: \passthrough{\lstinline!a.equals(b)!}
    \(\implies\) \passthrough{\lstinline!a.hashCode() == b.hashCode()!}.
  \item
    Proprietà desiderabile (riduzione collisioni):
    \passthrough{\lstinline"!a.equals(b)"} \(\implies\) preferibilmente
    \passthrough{\lstinline"a.hashCode() != b.hashCode()"}.
  \item
    Da Java 7 si usa l'utility di sistema:
    \passthrough{\lstinline!Objects.hash(campo1, campo2, ...)!}.
  \end{itemize}
\end{itemize}

\paragraph{\texorpdfstring{2. \texttt{TreeSet\textless{}E\textgreater{}}
(Slide
15-16)}{2. TreeSet\textless E\textgreater{} (Slide 15-16)}}\label{treesete-slide-15-16}

\begin{itemize}
\tightlist
\item
  \textbf{Struttura interna}: Albero binario di ricerca auto-bilanciante
  (\textbf{Red-Black Tree}).
\item
  \textbf{Complessità}: \passthrough{\lstinline!add!},
  \passthrough{\lstinline!remove!}, \passthrough{\lstinline!contains!}
  richiedono tempo logaritmico \textbf{\(O(\log n)\)}.
\item
  \textbf{Ordinamento garantito}: Gli elementi sono
  \textbf{costantemente mantenuti ordinati}:

  \begin{itemize}
  \tightlist
  \item
    Ordinamento naturale: richiede che gli elementi implementino
    \passthrough{\lstinline!Comparable<E>!}.
  \item
    Ordinamento esplicito: passando un
    \passthrough{\lstinline!Comparator<E>!} al costruttore:
    \passthrough{\lstinline!new TreeSet<>(comparator)!}.
  \end{itemize}
\item
  \textbf{Metodi speciali per insiemi ordinati}:

  \begin{itemize}
  \tightlist
  \item
    \passthrough{\lstinline!E first()!}: elemento minimo.
  \item
    \passthrough{\lstinline!E last()!}: elemento massimo.
  \item
    \passthrough{\lstinline!TreeSet<E> subSet(E fromElement, E toElement)!}:
    sottoinsieme nell'intervallo \([fromElement, toElement)\).
  \end{itemize}
\end{itemize}

\begin{center}\rule{0.5\linewidth}{0.5pt}\end{center}

\subsubsection{\texorpdfstring{1.5 Mappe Associative:
\texttt{Map\textless{}K,\ V\textgreater{}} (Slide
17-19)}{1.5 Mappe Associative: Map\textless K, V\textgreater{} (Slide 17-19)}}\label{mappe-associative-mapk-v-slide-17-19}

Una mappa modella una funzione matematica \(K \to V\) (tabella di
associazioni chiave-valore): * Le \textbf{chiavi (\(K\)) sono univoche}
(formano un \passthrough{\lstinline!Set<K>!}). * A ogni chiave è
associato esattamente un valore (\(V\)). I valori possono ripetersi. *
Metodi cardine: * \passthrough{\lstinline!V put(K key, V value)!}:
inserisce o rimpiazza l'associazione. *
\passthrough{\lstinline!V get(Object key)!}: restituisce il valore (o
\passthrough{\lstinline!null!} se la chiave non esiste). *
\passthrough{\lstinline!boolean containsKey(Object key)!} /
\passthrough{\lstinline!boolean containsValue(Object value)!} *
\passthrough{\lstinline!Set<K> keySet()!}: restituisce la vista
dell'insieme delle chiavi. *
\passthrough{\lstinline!Collection<V> values()!}: restituisce la
collezione dei valori. *
\passthrough{\lstinline!Set<Map.Entry<K, V>> entrySet()!}: restituisce
le coppie chiave-valore. *
\textbf{\passthrough{\lstinline!HashMap<K, V>!}}: implementazione
standard basata su hashing delle chiavi, con operazioni
\passthrough{\lstinline!put!} e \passthrough{\lstinline!get!} in
\textbf{\(O(1)\)}.

\begin{center}\rule{0.5\linewidth}{0.5pt}\end{center}

\subsubsection{\texorpdfstring{1.6 Iteratori e l'Errore
\texttt{ConcurrentModificationException} (Slide
21-23)}{1.6 Iteratori e l'Errore ConcurrentModificationException (Slide 21-23)}}\label{iteratori-e-lerrore-concurrentmodificationexception-slide-21-23}

Tutte le collezioni forniscono un \passthrough{\lstinline!Iterator<E>!}
(\passthrough{\lstinline!hasNext()!}, \passthrough{\lstinline!next()!},
\passthrough{\lstinline!remove()!}). Il ciclo
\passthrough{\lstinline!for-each!} di Java è uno zucchero sintattico
compilato tramite \passthrough{\lstinline!Iterator!}:

\begin{lstlisting}[language=Java]
for (Integer i : list) {
    System.out.println(i);
}
\end{lstlisting}

\paragraph{\texorpdfstring{La Trappola della Modifica Concorrente
(\emph{Fail-Fast Iterator}, Slide
23)}{La Trappola della Modifica Concorrente (Fail-Fast Iterator, Slide 23)}}\label{la-trappola-della-modifica-concorrente-fail-fast-iterator-slide-23}

\begin{lstlisting}[language=Java]
// CODICE ERRATO:
int count = 0;
for (Integer i : list) {
    System.out.println(i);
    if (count == 2)
        list.remove(count); // RUNTIME ERROR: java.util.ConcurrentModificationException!
    count++;
}
\end{lstlisting}

\begin{itemize}
\item
  \textbf{Perché crasha?}: Le collezioni Java mantengono internamente un
  contatore delle modifiche strutturali
  (\passthrough{\lstinline!modCount!}). Quando si invoca
  \passthrough{\lstinline!list.remove()!},
  \passthrough{\lstinline!modCount!} viene incrementato. Al ciclo
  successivo, la chiamata implicita a
  \passthrough{\lstinline!it.next()!} scopre che la collezione è
  cambiata alle sue spalle (\emph{fail-fast}) e lancia immediatamente
  \passthrough{\lstinline!ConcurrentModificationException!} per
  prevenire stati incoerenti.
\item
  \textbf{La Soluzione Corretta (Slide 23)}: Usare esplicitamente
  l'iteratore e invocare il \textbf{suo} metodo
  \passthrough{\lstinline!remove()!}, che aggiorna coerentemente lo
  stato interno dell'iteratore:

\begin{lstlisting}[language=Java]
for (Iterator<Integer> it = list.iterator(); it.hasNext(); ) {
    Integer i = it.next();
    if (count == 2)
        it.remove(); // PERFETTAMENTE LEGALE E SICURO!
    count++;
}
\end{lstlisting}
\end{itemize}

\begin{center}\rule{0.5\linewidth}{0.5pt}\end{center}

\subsubsection{\texorpdfstring{1.7 Algoritmi di Ordinamento:
\texttt{Comparable} vs \texttt{Comparator} (Slide
24-29)}{1.7 Algoritmi di Ordinamento: Comparable vs Comparator (Slide 24-29)}}\label{algoritmi-di-ordinamento-comparable-vs-comparator-slide-24-29}

La classe di utilità \passthrough{\lstinline!Collections!}
(\passthrough{\lstinline!java.util.Collections!}) fornisce due varianti
sovraccaricate del metodo statico \passthrough{\lstinline!sort!}:

\paragraph{Variante 1: Ordinamento
Naturale}\label{variante-1-ordinamento-naturale}

\begin{lstlisting}[language=Java]
public static <T extends Comparable<? super T>> void sort(List<T> list)
\end{lstlisting}

\begin{itemize}
\tightlist
\item
  \textbf{Perché \passthrough{\lstinline!<? super T>!} nella firma?
  (Slide 26 - Domanda Classica d'Esame)}: Supponiamo di avere
  \passthrough{\lstinline!class Student extends Person!}. Se solo
  \passthrough{\lstinline!Person!} implementa
  \passthrough{\lstinline!Comparable<Person>!}, allora
  \passthrough{\lstinline!Student!} eredita il confronto basato su
  \passthrough{\lstinline!Person!}. Se la firma fosse
  \passthrough{\lstinline!<T extends Comparable<T>>!}, allora
  \passthrough{\lstinline!Student!} dovrebbe implementare tassativamente
  \passthrough{\lstinline!Comparable<Student>!}. Con
  \passthrough{\lstinline!Comparable<? super T>!}, una lista di
  \passthrough{\lstinline!Student!} può essere ordinata legalmente
  usando il \passthrough{\lstinline!compareTo!} della superclasse
  \passthrough{\lstinline!Person!}
  (\passthrough{\lstinline!Comparable<Person>!}), garantendo il massimo
  riuso polimorfico!
\end{itemize}

\paragraph{\texorpdfstring{Variante 2: Ordinamento On-Demand con
\texttt{Comparator}}{Variante 2: Ordinamento On-Demand con Comparator}}\label{variante-2-ordinamento-on-demand-con-comparator}

\begin{lstlisting}[language=Java]
public static <T> void sort(List<T> list, Comparator<? super T> cmp)
\end{lstlisting}

\begin{itemize}
\tightlist
\item
  Permette di definire criteri di ordinamento arbitrari senza toccare la
  classe sorgente, passando:

  \begin{enumerate}
  \def\labelenumi{\arabic{enumi}.}
  \tightlist
  \item
    Una classe dedicata
    (\passthrough{\lstinline!class StudentComparator implements Comparator<Student>!}).
  \item
    Una \textbf{classe anonima}
    (\passthrough{\lstinline!new Comparator<Student>() \{ ... \}!}).
  \item
    Un'\textbf{espressione Lambda}
    (\passthrough{\lstinline!(s1, s2) -> s2.getMatricola() - s1.getMatricola()!}).
  \end{enumerate}
\end{itemize}

\begin{center}\rule{0.5\linewidth}{0.5pt}\end{center}

\subsection{\texorpdfstring{2. Riscontro Pratico nel Codice
(\texttt{Lezione15}/\texttt{Lezione16})}{2. Riscontro Pratico nel Codice (Lezione15/Lezione16)}}\label{riscontro-pratico-nel-codice-lezione15lezione16}

Come osservato durante l'ispezione dei sorgenti: 1. Nella Lezione 15 il
docente ha introdotto tutte le strutture viste nella Slide T16
all'interno di
\href{file:///home/wally/Documenti/GitHub/Programmazione-II/25-26/Teoria-e-Esercitazioni/Codice-20260902/Lezione15/MavenDate/src/main/java/it/oop/ui/MainDate.java}{MainDate.java}:
* \passthrough{\lstinline!List<Date> dateList = new ArrayList<>();!}
(con duplicato \passthrough{\lstinline!AmericanDate(7,11,2025)!}
\(\implies\) dimensione 4). *
\passthrough{\lstinline!Set<Date> dateSet = new HashSet<>();!}
(duplicato scartato \(\implies\) dimensione 3). *
\passthrough{\lstinline!Map<String, FormattedDateConverter> converter = new HashMap<>();!}
* Ordinamento naturale:
\passthrough{\lstinline!Collections.sort(dateList);!} * Ordinamento con
\passthrough{\lstinline!Comparator!} via classe anonima e lambda su
\passthrough{\lstinline!DateInterval!}. 2. Nella Lezione 16 il docente
compie il passo architetturale successivo: separa la gestione degli
errori creando il package \passthrough{\lstinline!it.oop.exception!} con
eccezioni personalizzate (\passthrough{\lstinline!IllegalDateException!}
e \passthrough{\lstinline!OrderdPairException!}), preparando il terreno
per la teoria della Slide T17.

\begin{center}\rule{0.5\linewidth}{0.5pt}\end{center}

Dimmi \textbf{``vai''} per procedere con il \textbf{Blocco 16} (Lezione
17 / Slide T17 --- \emph{Error Handling with Exceptions}, gerarchia
\passthrough{\lstinline!Throwable!}/\passthrough{\lstinline!Exception!}/\passthrough{\lstinline!RuntimeException!},
eccezioni controllate vs non controllate, costrutto
\passthrough{\lstinline!try-catch-finally!},
\passthrough{\lstinline!try-with-resources!} e il codice di
\passthrough{\lstinline!Lezione16/MavenDate!}).


\end{document}
